% =============================================================================
%  1-Introduccion.tex — Introducción / Planteamiento del problema
% =============================================================================

\section{Introducción: Contexto Socio-histórico}

También llamada \textit{Carta de los pieles rojas a las caras pálidas}, es la respuesta del
\textit{Gran jefe Seattle} enviada en 1854 al presidente estadounidense Franklin Pierce,
quien le propone al jefe nativo americano comprar sus tierras y reubicar a su gente en
una \textit{Reserva Ecológica}. La respuesta del nativo es contundente y propia de
la sabiduría de la tribu:

\textbf{¿Cómo se puede comprar o vender cielo, el calor de la tierra? La idea nos resulta
extraña.} \cite{perry1970home}.

Aunque la intención del \textit{Jefe Seattle} era hacer saber al nuevo gobierno de los Estados
Unidos de América la profunda desigualdad en ideas, el texto se convirtió en uno de los manifiestos
más profundos y citados en el ámbito de la ecología y la preservación de los espacios naturales.
Recordemos que los EUA experimentaron un rápido crecimiento y expansión hacia su territorio Oeste,
lo que se conoce popularmente como \textit{El Viejo y Salvaje Oeste}, territorio accidentado por
su orografía y aún inexplorado para el gobierno de las nacientes trece colonias del siglo
pasado (1776). El objetivo del mismo era adueñarse de los recursos naturales que había en esa
zona, al igual que expandir las vías ferroviarias para intercomunicar sus puertos de Este a Oeste.
Sin embargo, esa zona ya tenía dueño: las llamadas \textit{Tribus Nativas Americanas}. Fue en 1803
cuando se consolidó la venta de la Louisiana por el presidente Thomas Jefferson por 15 millones
de dólares \cite{turner1920frontier}, momento en que el proceso de exploración y conquista dio inicio. A pesar
de ello, no fue hasta 1830, con la \textbf{Ley de Traslado Forzoso de los Indios}, que inicia el
desplazamiento de los \textit{Indios} al oeste del Río Misisipi en el llamado \textit{``Sendero de las Lágrimas''}.

Con esta ideología de conquista, el gobierno estadounidense emprendió diferentes campañas
de índole política, económica y militar, entre las que destacan:
\begin{enumerate}
    \item \textbf{1848 — Tratado de Guadalupe Hidalgo:} Tras la guerra con México, este país cede más de la mitad de
        su territorio a EE. UU. (incluyendo California, Nevada, Utah y Nuevo México) \cite{white1991it}.

    \item \textbf{1848–1849 — Fiebre del Oro de California:} El descubrimiento de oro en Sutter's Mill desata
        una migración masiva. Los aventureros y mineros que llegaron en masa pasaron a ser conocidos
        como los ``Forty-Niners'' (los del 49) \cite{white1991it}.

    \item \textbf{1853 — Llegada de Levi Strauss y nacimiento de los Jeans:} El inmigrante bávaro Levi Strauss
        llega a San Francisco para abrir un negocio de telas y mercería durante la Fiebre del Oro.
        Al notar que los mineros necesitaban pantalones capaces de resistir el trabajo rudo, comenzó
        a fabricar prendas de lona reforzada. Más tarde, en 1873, junto al sastre Jacob Davis, patenta los
        icónicos pantalones de mezclilla con remaches de cobre (\textit{blue jeans}),
        transformando la vestimenta laboral y consolidando un imperio textil \cite{downey1992levis}.

    \item \textbf{1862 — Ley de Asentamientos Rurales (\textit{Homestead Act}):} El gobierno ofrece
        tierras baratas o gratuitas a los colonos en el Oeste, acelerando la invasión de territorios
        nativos \cite{white1991it, brown1970bury}.

    \item \textbf{1869 — El Primer Ferrocarril Transcontinental:} Se completa la línea
        ferroviaria de \textit{Union Pacific} y \textit{Central Pacific} en Promontory Summit, Utah.
        El comercio ferroviario conecta ambas costas del país, permitiendo el traslado masivo de mercancías,
        ganado y recursos naturales, pero fragmentando irreversiblemente las tierras de caza de los
        indígenas \cite{white1991it, turner1920frontier}.
\end{enumerate}
